\documentclass{csc2059practical}

\setmodulecode{CSC2059}
\setmoduletitle{Data Structures and Algorithms}
\setpracticalnumber{0}
\setpracticaltitle{CSC2059 Development Environment}
\setpracticalsubtitle{Pre-lab setup: Windows, macOS and Linux}
\setpracticaldate{Academic Year 2026--27}


\begin{document}

\maketitle

\section{Introduction}

Complete these checks before your first practical; installation is pre-lab preparation. On School computers, the software is already installed: follow the laboratory route below. On your own computer, follow only the section for your operating system, then complete the final checks.

\begin{infobox}
\textbf{What each tool does}
\begin{description}
\item[VS Code] Edits your files; its integrated terminal runs commands.
\item[Git] Records versions of your work; GitHub stores pushed commits online.
\item[C++ compiler] Translates C++ source into a program.
\item[Make] Runs the build and test instructions supplied in a Makefile.
\item[Tests] Check whether a program behaves as expected.
\end{description}
VS Code is the editor used here; it is a different product from Visual Studio.
\end{infobox}

\section{School Windows Computers}
Open VS Code and choose \textbf{Terminal $\rightarrow$ New Terminal}. Select \textbf{CSC2059 Terminal} from the terminal profile selector if another shell opens. Skip the personal-computer installation sections and go to Section~\ref{sec:checks}.

The terminal and VS Code's Git interface use two Git installations, configured to share your personal \mintinline{text}{I:\.gitconfig} file. Settings saved there persist across laboratory resets unless you delete that file from the root of your I: drive. Practical 1 sets your name, email, editor and Windows line-ending preference. Keep your practical repositories on I: too.

\section{Your Own Windows Computer}
Follow these Windows steps in order. They install the same tool families as the laboratory environment; your personal computer does not use the laboratory I: drive.

\subsection{Install Visual Studio Code}

Download Visual Studio Code for Windows from:

\begin{center}
\url{https://code.visualstudio.com/Download}
\end{center}

Run the installer and accept the normal installation options.

\subsection{Install Git for Windows}

Download Git for Windows from:

\begin{center}
\url{https://git-scm.com/download/win}
\end{center}

Run the installer. You may accept the default options during installation, with one important exception:

\begin{quote}
\textbf{When asked to choose Git's default editor, select Visual Studio Code.}
\end{quote}

This is one reason Visual Studio Code should be installed \emph{before} Git.

Git Bash will also be installed as part of Git for Windows. Although Git Bash may be useful for other purposes, \textbf{do not use Git Bash as the terminal for CSC2059 practical work}. Later in this guide you will configure the terminal used throughout the module.

\subsection{Install MSYS2}

Download the MSYS2 installer from:

\begin{center}
\url{https://www.msys2.org/}
\end{center}

When asked for the installation location, use:

\begin{minted}{text}
C:\Apps\MSYS64
\end{minted}

Complete the installation normally.

\subsection{Update MSYS2}

Open the \textbf{MSYS2 UCRT64} terminal from the Windows Start menu.

At the prompt, enter:

\begin{minted}{bash}
pacman -Suy
\end{minted}

Accept the updates when prompted.

During a core-system update, MSYS2 may tell you that all MSYS2 processes, including the current terminal, must be closed. \textbf{This is normal.} Allow the terminal to close.

If this happens, reopen the \textbf{MSYS2 UCRT64} terminal and run:

\begin{minted}{bash}
pacman -Suy
\end{minted}

again.

Continue until the update completes without asking you to restart the terminal.

\subsection{Install the CSC2059 Development Tools}

Make sure that you are using the \textbf{MSYS2 UCRT64} terminal.

Install the GCC C and C++ compilers:

\begin{minted}{bash}
pacman -S mingw-w64-ucrt-x86_64-gcc
\end{minted}

Install the GDB debugger:

\begin{minted}{bash}
pacman -S mingw-w64-ucrt-x86_64-gdb
\end{minted}

Install Make:

\begin{minted}{bash}
pacman -S make
\end{minted}

Finally, install Git within MSYS2:

\begin{minted}{bash}
pacman -S git
\end{minted}

Accept the default choices when \texttt{pacman} asks for confirmation.

\subsection{Configure the Git Prompt}

The Git prompt displays useful information about the current Git repository and branch while you work in the terminal.

Using \textbf{Windows File Explorer}, navigate to:

\begin{minted}{text}
C:\Program Files\Git\etc\profile.d
\end{minted}

Copy the file:

\begin{minted}{text}
git-prompt.sh
\end{minted}

Now navigate to:

\begin{minted}{text}
C:\Apps\MSYS64\etc\profile.d
\end{minted}

and paste \texttt{git-prompt.sh} into this folder.

\subsection{Create the CSC2059 Terminal in Visual Studio Code}

Open Visual Studio Code.

Open the Command Palette using \texttt{Ctrl+Shift+P} and search for:

\begin{quote}
\texttt{Preferences: Open User Settings (JSON)}
\end{quote}

This opens the Visual Studio Code \texttt{settings.json} file.

Add the following terminal profile to the Windows terminal profiles in your settings:

\begin{minted}{json}
"terminal.integrated.profiles.windows": {
    "CSC2059 Terminal": {
        "path": "C:\\Apps\\MSYS64\\usr\\bin\\bash.exe",
        "args": [
            "--login",
            "-i"
        ],
        "env": {
            "CHERE_INVOKING": "1",
            "MSYSTEM": "UCRT64"
        }
    }
},
"terminal.integrated.defaultProfile.windows": "CSC2059 Terminal"
\end{minted}

\textbf{Important:} If your \texttt{settings.json} file already contains a
Windows terminal profiles section, do not create a second one. Add the \texttt{CSC2059 Terminal} entry inside the existing section.

Save the file and restart Visual Studio Code.

\subsection{Share Your Personal Git Configuration}
Git for Windows and MSYS2 Git may otherwise use different home directories. Give both the same global configuration file:
\begin{enumerate}
\item In File Explorer's address bar, enter \texttt{\%USERPROFILE\%}. Copy the full folder path it opens (your Windows profile, not your GitHub username).
\item In Windows, search for \textbf{Edit environment variables for your account}. Add a \emph{user} variable. Its name is:
\begin{terminal}
GIT_CONFIG_GLOBAL
\end{terminal}
\item Set its value to that full path followed by \texttt{\textbackslash{}.gitconfig}, for example:
\begin{terminal}
C:\Users\Jane\.gitconfig
\end{terminal}
Use your actual path, without surrounding quotation marks. Keep any existing configuration file; do not replace its contents.
\item Close all VS Code and MSYS2 windows, then reopen them so they inherit the variable.
\end{enumerate}
Practical 1 will save your personal settings into this shared file. This is a one-time personal-Windows setup; the School image already supplies its own configuration location.

\subsection{Make the VS Code Launcher Available}
In the CSC2059 Terminal, run \texttt{code --version}. If it prints a version, skip this subsection. Otherwise, for a standard VS Code installation, copy this entire block into that terminal once:
\begin{bashcode}
cat >> ~/.bash_profile <<'EOF'

# CSC2059: add only the VS Code launcher directory.
for csc2059_code_bin in "/c/Program Files/Microsoft VS Code/bin" \
  "$(cygpath -u "$LOCALAPPDATA")/Programs/Microsoft VS Code/bin"
do
    if [ -f "$csc2059_code_bin/code" ]; then
        export PATH="$PATH:$csc2059_code_bin"
        break
    fi
done
unset csc2059_code_bin
EOF
\end{bashcode}
This adds a startup setting for the two usual VS Code installation locations. Close and reopen the terminal, then retry \texttt{code --version}. A custom installation location needs its own launcher path; ask for help if the check still fails. Do not import the entire Windows PATH to fix this one command.

\section{Your Own macOS Computer}
Install \href{https://code.visualstudio.com/docs/setup/mac}{Visual Studio Code for macOS}. Open the macOS Terminal and run:
\begin{bashcode}
xcode-select --install
\end{bashcode}
Follow the installer for Apple's Command Line Tools, which supply the compiler, Git and Make. If the tools are already installed, proceed to the checks below.

In VS Code, open the Command Palette (\texttt{Cmd+Shift+P}) and run \textbf{Shell Command: Install 'code' command in PATH}. Restart the terminal afterwards. Use VS Code's normal integrated terminal; do not install MSYS2 or the Windows CSC2059 terminal profile.

Apple's \texttt{g++} command can report \textbf{Apple Clang}. This is expected on this platform.

\section{Your Own Linux Computer}
Install VS Code using the \href{https://code.visualstudio.com/docs/setup/linux}{official Linux instructions} for your distribution. On \textbf{Ubuntu or Debian}, install the remaining tools with:
\begin{bashcode}
sudo apt update
sudo apt install build-essential git
\end{bashcode}
Other distributions need their equivalent C++ compiler, Make and Git packages. Use VS Code's normal integrated terminal; no MSYS2 installation or Windows terminal profile is needed.

\section{Check Your Setup}\label{sec:checks}
On every platform, open VS Code's integrated terminal and run:
\begin{bashcode}
git --version
g++ --version
make --version
code --version
\end{bashcode}
Each command should print version information, not a command-not-found error. Exact versions can differ between computers.

\textbf{Windows only:} select CSC2059 Terminal and also run:
\begin{bashcode}
which g++
which make
which git
echo $MSYSTEM
\end{bashcode}
The expected results, in order, are:
\begin{terminal}
/ucrt64/bin/g++
/usr/bin/make
/usr/bin/git
UCRT64
\end{terminal}
These paths do not apply to macOS or Linux. On Windows, \texttt{gdb --version} also checks the debugger installed above; debugging is not required for the first Git exercise.

\subsection{Editor Assistance}
In VS Code's Extensions view, install or check that \textbf{C/C++} by Microsoft is available. It provides code assistance, not the compiler itself. If prompted to select a compiler, use the compiler installed above (UCRT64 \texttt{g++.exe} on Windows, Apple's compiler on macOS, or your distribution's compiler on Linux).

\subsection{A Small Compile-and-Run Check}
Create a separate folder called \texttt{setup-check} using your file browser and open it with \textbf{File $\rightarrow$ Open Folder}. Create and save \texttt{setup-check.cpp} containing:
\begin{cppcode}
#include <iostream>
int main()
{
    std::cout << "CSC2059 setup works!\n";
}
\end{cppcode}
No prior C++ knowledge is needed for this check. Open a new terminal in this folder and compile it:
\begin{bashcode}
g++ setup-check.cpp -o setup-check
\end{bashcode}
On Windows, run:
\begin{bashcode}
./setup-check.exe
\end{bashcode}
On macOS or Linux, run:
\begin{bashcode}
./setup-check
\end{bashcode}
You should see \texttt{CSC2059 setup works!}. This simple command checks installation; later practicals use their supplied Makefiles and build options.

\section{Ready for Practical 1}
You are ready when the tool checks and small program work. Use CSC2059 Terminal on Windows and the normal VS Code integrated terminal on macOS or Linux. If a check fails, ask the teaching team for help before starting the practical. Practical 1 will configure your Git identity and editor.

\end{document}
